\documentclass{article}

\usepackage{amsmath,amssymb}
\usepackage{xurl}
\usepackage[hidelinks]{hyperref}
\setlength{\emergencystretch}{2em}

\input{command}

\title{The Square-Dimension Restriction in Wolf's Generic Normal Form\\
{\large (restriction resolved 2026-08-04)}}
\author{}
\date{2026-06-30, updated 2026-08-04}

\begin{document}
\maketitle
\gapnote{scope-restriction}{resolved}

\section{The assertion}

Proposition~2.9 of Ref.~\cite{Wolf2012Quantum} gives a normal form for
completely positive maps with full Kraus rank.\footnote{Local source:
\path{Notes/WolfNoteTexSource/ch02_representations.tex}, lines~923--929.}
Let
\begin{align}
    T&:\mathcal{M}_{d_1}\longrightarrow\mathcal{M}_{d_2}
    \label{eq:wolf_prop28_square_case_source_map}
\end{align}
be such a map. The proposition states that there are invertible completely
positive maps $\Phi_i\in\mathcal{B}(\mathcal{M}_{d_i})$, each of Kraus rank
one, such that
\begin{align}
    T'&=\Phi_2\circ T\circ\Phi_1,
    \label{eq:wolf_prop28_square_case_filtered_map}\\
    T'(\one)&\in\C\one,
    \label{eq:wolf_prop28_square_case_output_identity}\\
    T'^{*}(\one)&\in\C\one.
    \label{eq:wolf_prop28_square_case_adjoint_identity}
\end{align}
Conditions~(\ref{eq:wolf_prop28_square_case_output_identity}) and
(\ref{eq:wolf_prop28_square_case_adjoint_identity}) express double
stochasticity in the source's sense. The input filtering $\Phi_1$ acts on
$\mathcal{M}_{d_1}$, and the output filtering $\Phi_2$ acts on
$\mathcal{M}_{d_2}$. Proposition~2.9 of Ref.~\cite{Wolf2012Quantum} imposes
no relation between $d_1$ and $d_2$.

\section{The point at issue}

The original formal theorem covered only the equal-dimension specialization
\begin{align}
    d_1&=d_2=D,
    \label{eq:wolf_prop28_square_case_equal_dimensions}\\
    T&:\MN{D}\longrightarrow\MN{D}.
    \label{eq:wolf_prop28_square_case_square_map}
\end{align}
Both filterings in this theorem act on $\MN{D}$.\footnote{Formal declaration:
\leanid{Wolf.exists\_normal\_form\_generic} in
\doclink{QICLean/Channel/LorentzNormalForm}.}

For the square map in
(\ref{eq:wolf_prop28_square_case_square_map}), the Choi matrix is
\begin{align}
    \tau
        &=(T\otimes\operatorname{id})
          (\ketbra{\Omega}{\Omega})
          \in\MN{D}\otimes\MN{D}.
    \label{eq:wolf_prop28_square_case_square_choi}
\end{align}
The original minimization therefore ranged over
$\operatorname{SL}(D,\C)\times\operatorname{SL}(D,\C)$. Its compactness and
trace--determinant equality arguments were also stated only for the square
space in (\ref{eq:wolf_prop28_square_case_square_choi}).\footnote{Formal
declarations:
\leanid{Wolf.infimum\_is\_attained} and
\leanid{Matrix.PosSemidef.pow\_card\_mul\_det\_re\_eq\_trace\_re\_pow\_iff}.}

At $D=2$, this square theorem supplies the full-Kraus-rank minimization result
used in the diagonal case of the qubit Lorentz normal form. It does not prove
the existence theorem combining the diagonal, non-diagonal, and singular
cases. The scalar-normalization and Lorentz-orbit gap in that theorem remains
recorded in
\path{docs/paper-gaps/wolf_prop2_11_lorentz_scalar_filtering_gap.tex}.

\section{The restricted statement}

Before the rectangular generalization, the formal development proved
\begin{align}
    T&:\MN{D}\longrightarrow\MN{D},
    \label{eq:wolf_prop28_square_case_restricted_domain}\\
    \tau=\operatorname{choi}(T)&\succ0,
    \label{eq:wolf_prop28_square_case_positive_choi}\\
    \tau\succ0
        &\Longrightarrow
          \exists\,\Phi_1,\Phi_2:\,
          \Phi_2\circ T\circ\Phi_1
          \text{ is doubly stochastic}.
    \label{eq:wolf_prop28_square_case_restricted_conclusion}
\end{align}
Here $\Phi_1$ and $\Phi_2$ are determinant-one filterings on $\MN{D}$. This
is the specialization $d_1=d_2$ of Proposition~2.9 of
Ref.~\cite{Wolf2012Quantum}. It is correct but narrower than the source's
rectangular statement.

\section{Resolution}

The rectangular generalization was completed on 2026-08-04 under
\href{https://github.com/LionSR/TNLean/issues/3616}{TNLean issue \#3616}.
It consists of two changes.
\begin{itemize}
    \item The declaration \leanid{Wolf.infimum\_is\_attained} was generalized
    in place to positive-definite operators on
    $\C^{d_2}\otimes\C^{d_1}$ and filterings in
    $\operatorname{SL}(d_1,\C)\times\operatorname{SL}(d_2,\C)$. The
    Frobenius-norm confinement now uses separate dimension constants for the
    two factors. The square caller remains unchanged because its dimensions
    unify.

    \item The rectangular proposition is formalized as
    \leanid{Wolf.exists\_normal\_form\_generic\_rect}. It applies to a
    completely positive map
    \begin{align}
        T&:\mathcal{M}_{d_1}\longrightarrow\mathcal{M}_{d_2}
        \label{eq:wolf_prop28_square_case_rectangular_map}
    \end{align}
    in the rectangular Kraus sense \leanid{IsKrausCP}, assuming that its
    rectangular Choi matrix \leanid{ChoiRectangular.choiMatrix} is positive
    definite. It produces a Kraus-rank-one filtering $\Phi_1$ on
    $\mathcal{M}_{d_1}$ and a Kraus-rank-one filtering $\Phi_2$ on
    $\mathcal{M}_{d_2}$, each represented by a determinant-one matrix, such
    that
    \begin{align}
        T'&=\Phi_2\circ T\circ\Phi_1,
        \label{eq:wolf_prop28_square_case_rectangular_filtered_map}\\
        T'(\one)&\in\C\one,
        \label{eq:wolf_prop28_square_case_rectangular_output_identity}\\
        T'^{*}(\one)&\in\C\one.
        \label{eq:wolf_prop28_square_case_rectangular_adjoint_identity}
    \end{align}
    These conditions are represented by
    \leanid{Wolf.DoublyStochasticRect}, using the trace-pairing adjoint.
\end{itemize}

\section{Verdict}

The scope restriction is resolved. Proposition~2.9 of
Ref.~\cite{Wolf2012Quantum} is now formalized with independent input and
output dimensions. The declaration
\leanid{Wolf.exists\_normal\_form\_generic} remains the unchanged
equal-dimension specialization. The remaining gap concerns Proposition~2.11
of Ref.~\cite{Wolf2012Quantum}, namely the qubit Lorentz classification, and
is recorded in
\path{docs/paper-gaps/wolf_prop2_11_lorentz_scalar_filtering_gap.tex}.

\bibliographystyle{alpha}
\bibliography{references}

\end{document}
